\begin{lemma}[Hidden-coordinate instance factorization]\label{lem:step-008-instance-factorization}
Under Proposition~\ref{prop:step-004-finite-hard-prior} and the
ideal experiment above, if \(J\sim\operatorname{Unif}[k]\) and, conditional
on each \(J=j\), the coordinate instances
\(\Xi_1,\ldots,\Xi_k\) are mutually independent with their common law
\(\mu_{N,M}\), then

\[
\mathcal L(\boldsymbol\Xi\mid J=j)=\mu_{N,M}^{\otimes k}
\quad\text{for every }j\in[k],
\tag{5}
\]

and consequently (2) holds. This remains true when \(\mu_{N,M}\) has
singleton support or repeated coordinate draws occur.
\end{lemma}

\begin{proof}
Let \(\mathcal Z=\operatorname{supp}(\mu_{N,M})\), which is a nonempty
finite set because Proposition~\ref{prop:step-004-finite-hard-prior}
supplies a finitely supported
probability law. For every \(j\in[k]\) and
\(\boldsymbol\xi=(\xi_1,\ldots,\xi_k)\in\mathcal Z^k\), the specified
conditional independence gives

\[
\begin{aligned}
\Pr\{J=j,\boldsymbol\Xi=\boldsymbol\xi\}
&=\Pr\{J=j\}
  \Pr\{\boldsymbol\Xi=\boldsymbol\xi\mid J=j\}\\
&=\frac1k\prod_{i=1}^k\mu_{N,M}(\{\xi_i\}).
\end{aligned}
\tag{6}
\]

The product on the right does not depend on \(j\). Summing (6) over
\(j\) gives

\[
\Pr\{\boldsymbol\Xi=\boldsymbol\xi\}
=\prod_{i=1}^k\mu_{N,M}(\{\xi_i\}),
\]

which is the product law in (2). Equation (6) then factors as

\[
\Pr\{J=j,\boldsymbol\Xi=\boldsymbol\xi\}
=\Pr\{J=j\}\Pr\{\boldsymbol\Xi=\boldsymbol\xi\},
\]

proving independence. If the prior has singleton support, all coordinates
are the same deterministic pair almost surely, but the same factorization
still proves both the product law and independence from \(J\). Calling one
coordinate hidden therefore creates no selection bias. \end{proof}

\begin{lemma}[Conditional i.i.d. tagged-product sample law]\label{lem:step-008-ideal-sample-law}
Under the primitive tagged-product definitions and
Lemma~\ref{lem:step-008-instance-factorization}, suppose that after
\((J,\boldsymbol\Xi)\) is generated, the tags
\(I_1,\ldots,I_n\) are drawn independently from
\(\operatorname{Unif}[k]\), independently of the preceding variables;
conditional on \((\boldsymbol\Xi,I_1,\ldots,I_n)\), the features are
independent with \(X_\ell\sim Q_{I_\ell}\); and
\(Y_\ell=\tau_{T_{I_\ell}}(X_\ell)\). Then for every realized
\(\boldsymbol\xi=((t_i,Q_i))_{i=1}^k\in\mathcal Z^k\) and every
\(j\in[k]\),

\[
\mathcal L\!\left(S^{\mathrm{id}}\mid
  \boldsymbol\Xi=\boldsymbol\xi,J=j\right)
=\left(P_{\boldsymbol Q}^{c_{\boldsymbol t}}\right)^n.
\tag{7}
\]

In particular, the conditional row kernel is independent of \(j\), and (3)
holds.
\end{lemma}

\begin{proof}
Fix \(\boldsymbol\xi=((t_i,Q_i))_{i=1}^k\) and \(j\). For any row index
\(\ell\), any \(i\in[k]\), \(x\in[N]\), and \(y\in\{0,1\}\), the tag,
feature, and deterministic-label rules give

\[
\begin{aligned}
&\Pr\left\{((I_\ell,X_\ell),Y_\ell)=((i,x),y)
  \mid\boldsymbol\Xi=\boldsymbol\xi,J=j\right\}\\
&\qquad=\frac1k Q_i(x)\,
  \mathbf1\{y=\tau_{t_i}(x)\}\\
&\qquad=P_{\boldsymbol Q}(i,x)\,
  \mathbf1\{y=c_{\boldsymbol t}(i,x)\}\\
&\qquad=P_{\boldsymbol Q}^{c_{\boldsymbol t}}
  \bigl(\{((i,x),y)\}\bigr).
\end{aligned}
\tag{8}
\]

The first line on the right uses neither \(j\) nor any property of the
learner. The pairs \((I_\ell,X_\ell)\) are mutually independent conditional
on \(\boldsymbol\Xi\): tags are drawn independently, and the conditional
feature draws use independent randomness across rows. Labels are
deterministic rowwise functions. Therefore the \(n\) labeled rows are
conditionally independent and all have the common law in (8), which proves
(7).

By Lemma~\ref{lem:step-008-instance-factorization}, the conditional law of
\(\boldsymbol\Xi\) itself is the same for every \(j\). Since (7) is also the
same for every \(j\), summing over \(J\) yields (3). Importantly, (3) is a
conditional i.i.d. statement for each fixed product instance. Unconditionally,
the rows share the random vector \(\boldsymbol\Xi\), so no stronger
unconditional i.i.d. assertion is made or needed.

Equation (8) remains valid for \(t_i=1\) or \(t_i=N+1\), when the labels in
that block are constant, and when \(Q_i\) is a point mass or has zeros. No
division by \(Q_i(x)\) occurs. It also defines an exact size-\(n\) dataset
when no sampled tag equals the later-designated hidden coordinate \(J\).
\end{proof}

\begin{lemma}[Persistence of selector independence through an asymmetric
learner]\label{lem:step-008-selector-independence}
Under Lemmas~\ref{lem:step-008-instance-factorization} and
\ref{lem:step-008-ideal-sample-law}, if \(A\) is any randomized map with the
domain and codomain in Section~\ref{sec:setup} and is run on
\(S^{\mathrm{id}}\) with fresh internal randomness, then

\[
J\ \perp\!\!\!\perp\
(\boldsymbol\Xi,S^{\mathrm{id}},H).
\tag{9}
\]

Equivalently, for every \(i\in[k]\),

\[
\Pr\{J=i\mid\boldsymbol\Xi,S^{\mathrm{id}},H\}=\frac1k
\quad\text{almost surely}.
\tag{10}
\]

No permutation invariance, tag symmetry, properness, or determinism of
\(A\) is required.
\end{lemma}

\begin{proof}
Because all involved spaces are finite, represent the randomized learner by
its stochastic kernel

\[
K_A(h\mid s):=\Pr\{A(s)=h\},
\qquad h\in\mathcal H_{k,N}.
\]

For a fixed instance vector \(\boldsymbol\xi\), let
\(p_{\boldsymbol\xi}(s)\) denote the product-sample mass furnished by
Lemma \ref{lem:step-008-ideal-sample-law}. Equations (6)--(8), followed by
the independent learner call, give for every atom with
\(\boldsymbol\xi\in\mathcal Z^k\),

\[
\begin{aligned}
&\Pr\{J=j,\boldsymbol\Xi=\boldsymbol\xi,
       S^{\mathrm{id}}=s,H=h\}\\
&\qquad=\frac1k
  \left(\prod_{i=1}^k\mu_{N,M}(\{\xi_i\})\right)
  p_{\boldsymbol\xi}(s)K_A(h\mid s).
\end{aligned}
\tag{11}
\]

Every factor after \(1/k\) is independent of \(j\). Summing (11) over
\(j\) shows that it equals

\[
\Pr\{J=j\}\Pr\{\boldsymbol\Xi=\boldsymbol\xi,
S^{\mathrm{id}}=s,H=h\},
\]

which proves (9), including zero-mass atoms by omission. The finite
conditional-probability formula gives (10) on every positive-mass atom.

The kernel \(K_A(h\mid s)\) may depend arbitrarily on every tag and every
row of \(s\). Thus (11) does not assert that the law of \(H\) is invariant
under permutations of tags. It uses only that \(J\) is an analysis-side
variable not supplied as an argument to \(A\), and that the sample kernel
in (7) does not depend on \(J\). The learner may also have the public prior
hardwired into its code; this changes \(K_A\) but introduces no \(j\)-factor.
\end{proof}

\begin{proposition}[Deterministic tagged-product risk decomposition]\label{prop:step-008-risk-decomposition}
Under the primitive product-distribution definition and Proposition~\ref{prop:step-007-risk-identity}, for every fixed
\(\boldsymbol t\in[N+1]^k\), every vector of probability distributions
\(\boldsymbol Q=(Q_1,\ldots,Q_k)\) on \([N]\), and every arbitrary
\(h\in\mathcal H_{k,N}\),

\[
\begin{aligned}
R_{P_{\boldsymbol Q}}(h,c_{\boldsymbol t})
&=\frac1k\sum_{i=1}^k e_i(h;\boldsymbol t,\boldsymbol Q)\\
&=\frac1k\sum_{i=1}^k
  R_{Q_i}(D_i h,\tau_{t_i}).
\end{aligned}
\tag{12}
\]

The equality is deterministic and does not require any distributional
symmetry of \(\boldsymbol Q\) or any structural property of \(h\).
\end{proposition}

\begin{proof}
Expanding the population risk on the finite tagged domain and using
\[
P_{\boldsymbol Q}(i,x)=k^{-1}Q_i(x)
\]
gives

\[
\begin{aligned}
R_{P_{\boldsymbol Q}}(h,c_{\boldsymbol t})
&=\sum_{i=1}^k\sum_{x\in[N]}
  \frac1k Q_i(x)
  \mathbf1\{h(i,x)\ne c_{\boldsymbol t}(i,x)\}\\
&=\frac1k\sum_{i=1}^k\sum_{x\in[N]}Q_i(x)
  \mathbf1\{h(i,x)\ne\tau_{t_i}(x)\}\\
&=\frac1k\sum_{i=1}^k e_i(h;\boldsymbol t,\boldsymbol Q).
\end{aligned}
\tag{13}
\]

Proposition~\ref{prop:step-007-risk-identity} identifies the
\(i\)-th inner sum with
\(R_{Q_i}(D_i h,\tau_{t_i})\), proving the second equality in (12). The
calculation remains exact for unequal block distributions, endpoint
thresholds, point masses, and arbitrary nonmonotone or tag-asymmetric
hypotheses. \end{proof}

\begin{proposition}[Selected-coordinate averaging for arbitrary randomized
learners]\label{prop:step-008-selected-risk-identity}
Under Proposition~\ref{prop:step-007-risk-identity},
Lemma~\ref{lem:step-008-selector-independence}, and
Proposition~\ref{prop:step-008-risk-decomposition}, for the ideal experiment
and every randomized map \(A\) specified above,

\[
\mathbb E\!\left[
e_J(H;\boldsymbol T,\boldsymbol Q)
\mid\boldsymbol\Xi,S^{\mathrm{id}},H
\right]
=R_{P_{\boldsymbol Q}}(H,c_{\boldsymbol T})
\quad\text{almost surely}.
\tag{14}
\]

Consequently,

\[
\begin{aligned}
\mathbb E\,
R_{Q_J}(D_JH,\tau_{T_J})
&=\mathbb E\,e_J(H;\boldsymbol T,\boldsymbol Q)\\
&=\mathbb E\,R_{P_{\boldsymbol Q}}(H,c_{\boldsymbol T}),
\end{aligned}
\tag{15}
\]

where the expectations retain all prior, sample, selector, and learner
randomness.
\end{proposition}

\begin{proof}
Conditional on \((\boldsymbol\Xi,S^{\mathrm{id}},H)\), every quantity
\(e_i(H;\boldsymbol T,\boldsymbol Q)\) is fixed. By (10), the remaining
coordinate selector is uniform, so the elementary finite
conditional-expectation identity restated above gives

\[
\begin{aligned}
\mathbb E\!\left[
e_J(H;\boldsymbol T,\boldsymbol Q)
\mid\boldsymbol\Xi,S^{\mathrm{id}},H
\right]
&=\sum_{i=1}^k\frac1k
  e_i(H;\boldsymbol T,\boldsymbol Q)\\
&=R_{P_{\boldsymbol Q}}(H,c_{\boldsymbol T}),
\end{aligned}
\]

where the last equality is the pathwise application of
Proposition~\ref{prop:step-008-risk-decomposition}. Taking expectations and
using the total-expectation part of the same restated finite identity proves
the second equality in (15), which is (4). The first equality in (15) is the
pathwise Proposition~\ref{prop:step-007-risk-identity} at the realized
\((J,H,T_J,Q_J)\).

The proof averages over an independent uniform \(J\) after conditioning on
the complete realized vector, sample, and output. It never replaces
\(e_J\) by \(e_1\), never permutes \(A\), and never assumes the errors
\(e_1,\ldots,e_k\) have the same distribution. Thus deliberate tag
asymmetry of \(A\) is fully covered. For \(k=2\) and \(k=3\), the
conditional sums are exactly \((e_1+e_2)/2\) and
\((e_1+e_2+e_3)/3\), respectively. Samples in which the selected tag does
not occur require no special conditioning and remain part of the same
unconditioned identity. \end{proof}
